
The authors have developed an impressively broad and flexible approach for regressing non-Euclidean object data on multivariate Euclidean predictors.
In theory, the only requirement is that the objects reside in a metric space, i.e., that a suitable distance is defined for them.
This distance is employed in a first-stage manifold learning of the object-valued responses, to embed them in a lower-dimensional latent space that aims to preserve their geodesic distances.
The authors then employ non-parametric regression using modern deep learning to regress the latent embeddings on the multivariate covariates.
Again, this non-parametric regression step is very general, and allows them to flexibly capture non-linearities and interaction effects without imposing parametric assumptions.
The fitted values from the latent non-parametric regression are subsequently used as low-dimensional predictors in a local \Frechet regression step, thereby leveraging the structure learned in the latent non-parametric regression with the \Frechet mapping using the appropriate distance metric and ensuring that the objects' geometric constraints are preserved in all predictions.
The generality of the proposed methodology means that it can be applied, as the authors demonstrate in their examples, to connectivity matrices and to univariate probability distributions alike without any modifications, and without making any restrictive assumptions on the form of covariate effects.
To our best understanding, the method is also readily applicable to, for example, spherical and compositional data.

The theoretical results provide a solid grounding for the framework and are one of the work's strengths.
The authors establish convergence rates for DFR that account for the error in both the modeling of the latent embeddings and in the local \Frechet regression step.
The data application, where DFR is used to model the effect of meteorological variables, days of the week, and aggregated service-related covariates on New York City taxi networks is another strength of the work.
This interesting example highlights the practical applicability of the framework to modern non-Euclidean data.
